Los experimentos de esta tarea permitieron comprobar en la práctica cómo escalan los diferentes enfoques algorítmicos frente al problema ``AniMaraton''. Quedó en evidencia que Fuerza Bruta es inviable para instancias grandes, lo que se demostró al tener que detener la ejecución manualmente en $N=100$ debido a los tiempos excesivos provocados por su naturaleza exponencial. Por otro lado, la Programación Dinámica demostró ser la mejor estrategia para asegurar la solución óptima global en tiempos estables, pagando un costo fijo e invariable de memoria de 21.6 MB dictado por su matriz de estados ($O(M \times E)$). Finalmente, las heurísticas Greedy validaron su utilidad como una alternativa extremadamente rápida y ligera si se tolera perder la optimalidad, destacando el rendimiento superior de Greedy 1 sobre Greedy 2 al incluir ambas restricciones (minutos y energía) dentro de su criterio local de selección.